\section{Contexto}
\label{sec:contexto}

\subsection{Integración de P1--P5}

A lo largo de la asignatura se han implementado cinco prácticas
independientes pero complementarias, cada una orientada a un eje
distinto del aprendizaje automático aplicado al dominio financiero
personal:

\begin{itemize}[leftmargin=1.5em,itemsep=0.2em]
  \item \textbf{P1 --- Predicción temporal.} Modelos de regresión y
        series temporales (Random Forest, HistGradientBoosting y
        ARIMA) para anticipar el gasto del mes siguiente.
  \item \textbf{P2 --- Clasificación multietiqueta.} Asignación de la
        categoría (\textit{Area}) a una transacción a partir de su
        descripción textual mediante un \texttt{SGDClassifier} sobre
        char $n$-gramas.
  \item \textbf{P3 --- OCR financiero.} Extracción del importe total
        de facturas reales mediante \texttt{PaddleOCR} y un scorer
        basado en Gradient Boosting.
  \item \textbf{P4 --- Agente conversacional.} Primer contacto con
        \texttt{LangGraph} y la orquestación de herramientas mediante
        LLM.
  \item \textbf{P5 --- Biometría facial y detección de anomalías.}
        Pipeline MTCNN + FaceNet + DenseNet201 para autenticación
        biométrica con \textit{liveness}, e \texttt{IsolationForest}
        para anomalías financieras.
\end{itemize}

La práctica final exige evolucionar la práctica 6, que integra las cinco anteriores bajo un único sistema
multiagente que conserva la fiabilidad de las soluciones individuales
ganando cohesión, persistencia multiusuario y observabilidad
profesional. El enunciado considera P4 cubierto, por lo
que la \textit{evolución incremental} (Sección~\ref{subsec:evoluciones})
recae sobre P2, P3 y P5.

\subsection{Requisitos funcionales}

La especificación de la práctica final incluye un conjunto de
requisitos transversales que el sistema debe satisfacer. La
Tabla~\ref{tab:requisitos} los enumera y los enlaza con la sección de
la propuesta en la que se justifica su cumplimiento.

\begin{table}[!htbp]
  \centering
  \footnotesize
  \begin{tabularx}{\linewidth}{@{}p{0.42\linewidth} X@{}}
    \toprule
    \textbf{Requisito} & \textbf{Sección de la propuesta} \\
    \midrule
    Histórico conversacional para el LLM &
      \S\ref{subsec:memoria-conversacional} \\
    Recuperación de chats previos &
      \S\ref{subsec:memoria-conversacional} \\
    Diagramas visuales dinámicos (XAI) &
      \S\ref{subsec:xai-visualizacion} \\
    Monitorización del sistema (MonitorAgent) &
      \S\ref{subsec:monitorizacion} \\
    Al menos dos roles diferenciados en LangGraph &
      \S\ref{subsec:roles} \\
    Módulos P1--P5 distribuidos vía API REST &
      \S\ref{subsec:microservicios} \\
    Diálogo natural con detección de intención &
      \S\ref{subsec:intencion} \\
    Uso de Langfuse &
      \S\ref{subsec:langfuse} \\
    Prompt engineering (ASPECCT) &
      \S\ref{subsec:aspect} \\
    Tres evoluciones incrementales (P2, P3, P5) &
      \S\ref{subsec:evoluciones} \\
    Evaluaciones cuantitativas &
      \S\ref{subsec:evoluciones} \\
    Pruebas unitarias y de estrés &
      \S\ref{subsec:tests} \\
    \bottomrule
  \end{tabularx}
  \caption{Mapeo entre los requisitos del enunciado y las secciones de la
  propuesta.}
  \label{tab:requisitos}
\end{table}

\subsection{Restricciones técnicas y de despliegue}

El sistema se ha diseñado con tres restricciones operativas
relevantes que han condicionado tanto la elección de tecnologías
como la arquitectura final.

\paragraph{Stack gratuito o de coste mínimo.} El LLM principal corre
en \textit{Groq Cloud} (plan gratuito con \textit{rate limits} de
30 \textit{requests}/min para \texttt{llama-3.3-70b-versatile}), la
base de datos vive en \textit{Supabase} (Postgres gestionado con
500\,MB en el plan \textit{free}), la observabilidad en
\textit{Langfuse Cloud} (50.000 \textit{traces}/mes gratuitos) y el
frontend en \textit{Vercel}. El \textit{backend} se sirve localmente
con \texttt{uvicorn} y se expone públicamente vía un túnel
\textit{ngrok} con dominio estático, ya que el plan gratuito de
Vercel no permite contenedores con dependencias nativas pesadas
(PaddlePaddle, PyTorch, OpenCV) en sus \textit{serverless functions}.

\paragraph{Compatibilidad de versiones.} \texttt{PaddleOCR 2.6.2}
exige \texttt{numpy<2.0}, lo que ancla todo el \textit{stack} a
\texttt{NumPy 1.26.4}. Esta restricción ha condicionado la elección del
resto de dependencias: \texttt{scikit-learn 1.5.0},
\texttt{torch 2.5.1}, \texttt{langgraph 1.x},
\texttt{facenet-pytorch} (en lugar de \texttt{deepface}, que arrastra
TensorFlow). El \textit{trade-off} aceptado es no poder usar
\texttt{numpy 2.x} ni las versiones más recientes de
\texttt{scikit-learn}; a cambio, todo el \textit{stack} convive en
un único \textit{entry point} (\texttt{uvicorn}) sin necesidad de
aislar P3 en un microservicio aparte.

\paragraph{Activos pesados con Git LFS.} Los modelos preentrenados
(\texttt{models/area\_classifier.joblib},
\texttt{models/ocr\_total\_extractor.joblib} y
\texttt{models/liveness\_kaggle.pth}, $\sim$78\,MB en total) viajan
por \textit{Git Large File Storage} para no inflar el historial del
repositorio. Sin habilitar \texttt{git lfs install} antes del
\texttt{clone}, los modelos se descargan como punteros de 134 bytes
y el \textit{backend} falla al cargar el clasificador de
\texttt{Area} en el primer turno. Adicionalmente, las dependencias
DL (PaddleOCR, PyTorch, MediaPipe) ocupan $\sim$2{,}5\,GB en
\texttt{.venv}, lo que descarta la idea inicial de empaquetar la
aplicación en una imagen Docker desplegable en plataformas
\textit{serverless} sin almacenamiento persistente.

\paragraph{Cumplimiento RGPD y aislamiento multiusuario.} El sistema
maneja datos personales sensibles (transacciones financieras y
\textit{embeddings} biométricos), por lo que el cumplimiento del
Reglamento General de Protección de Datos no es una característica
opcional sino una restricción de diseño. Esto se traduce en cuatro
exigencias concretas: (i) los \textit{embeddings} faciales se
almacenan cifrados en disco con \texttt{Fernet} (AES-128) sobre una
clave derivada con \texttt{PBKDF2-HMAC-SHA256} a 310.000 iteraciones,
nunca en claro; (ii) las contraseñas se almacenan con \texttt{bcrypt}
(\textit{cost factor} 12) sin recurrir a \texttt{passlib} ---
incompatible con \texttt{bcrypt 4.x}; (iii) las notificaciones a
Telegram nunca exponen importes ni descripciones salvo que el
usuario active \texttt{notification\_level='full'}; y (iv) la
detección de anomalías nunca compara las transacciones de un usuario
contra el histórico de otro.
